% filepath: PID-in-Modern-Cherenkov/Sections/conclusion.tex
% Content for the "Conclusion" subsection

This review introduced the field of Cherenkov imaging detectors, and covered modern developments towards 
precision PID in upcoming high-energy physics experiments. The detectors of today will be able to further our 
understanding of quantum chromodynamics, gluons in hadron events, and elusive high-energy neutrinos. 

Highlights from the developments into the detectors include the push to improve spatial resolution of photosensors. There are also the improvements in timing resolution, with HRPPDs reaching an 
impressive 15ps single photon time resolution. The lifetime and radiation resistance of the photosensors has also been developed, with novel machining techniques such 
as ALD. Furthermore, ML for particle classification is one of the fastest moving areas of the field. They have minimal 
computation time compared to other statistical PID methods depending on the model size, and have been shown to match the 
accuracy of the old methods with large detailed models.

There is no shortage of improvements to be made to these detectors. As physics programmes demand higher luminosities and
higher energies, the technology must adapt. There are clear limitations in the high-rate capability of photosensors and 
their DAQ systems. The size of data per second expected to be generated by near future detectors is immense. The cost
of scientific programmes is ever more scrutinised, and so refinements towards the manufacturing of detector components
are an area of particular focus.

Nevertheless, it is an exciting time for detector physicists; the FAIR international accelerator facility and the HL LHC 
upgrades are both due to complete construction before 2035. The Electron-Ion Collider at BNL is due to begin operations
after 2035. Hyper-K, the successor to Super-K in Japan, is a neutrino detector that began civil engineering works and excavation 2 years
ago. 

Cherenkov imaging detectors have and will continue to be instrumental in high-energy physics research. They are a tried 
and true method of identifying charged particles. In the modern era, the detectors are beginning to move beyond auxiliary 
devices for separate mass hypotheses. Rather, they are playing vital roles in entire event reconstruction, allowing physicists 
to better study the interactions of our universe.
